\section{API Development Workflow}\label{apdx:api_construction}
In this section, we detail the methodology for leveraging LLMs to automate API construction. We propose a semi-automated workflow wherein users need only supply exemplar tasks performed within the target application; the LLM then autonomously generates both the necessary API code and corresponding test cases. The workflow comprises three primary stages: requirement analysis, API implementation, and test case generation.

\vpara{Requirement Analysis}
During the requirements analysis phase, users provide a set of task examples related to the target application as input. The workflow leverages the LLM to analyze these task instances, extracting the essential functionalities required for task completion. It then compares these requirements against the existing API interface definitions to identify potential gaps. If uncovered functionalities are detected, the system automatically generates new API interfaces along with their corresponding parameter specifications.

Notably, we limit the generated interfaces to encapsulate only general-purpose functionalities, thereby avoiding excessive complexity and the proliferation of APIs. This design choice mitigates implementation difficulty and reduces the adaptation burden on the agent.

\vpara{API Implementation}
Upon obtaining the interface definitions, the workflow systematically iterates over each interface and its associated parameters. For each specification, it leverages the designated Python libraries of the target application to implement the corresponding API functionalities. Additionally, the workflow incorporates error-handling mechanisms and logging to facilitate human debugging and maintenance. This automated approach not only streamlines API development but also enhances consistency and reusability across different application contexts.

\vpara{Test Case Generation}
Following the implementation of API functionalities, the workflow conducts fundamental unit testing to ensure the correctness and robustness of each API. Specifically, the testing process verifies: (1) whether the API can be invoked without runtime errors, and (2) whether the API returns correct results across a range of parameter inputs. For API implementations that fail these tests, the workflow provides detailed error feedback to the API implementation module, which then autonomously attempts corrections until the APIs pass all tests.

This automated methodology substantially lowers the manual effort required, enabling the creation of application-specific API sets with minimal human intervention. As a result, the barrier for users to develop APIs for diverse applications is significantly reduced. We have developed API sets for multiple commonly used default applications in Ubuntu and integrated them into our Ubuntu virtual machine environment.

\newpage